\documentclass[11pt, twocolumn, landscape, a4paper]{article}

\usepackage[margin=1cm]{geometry}
\usepackage[utf8]{inputenc}
\usepackage[spanish]{babel}
\usepackage{amssymb, amsmath, amsbsy}
\usepackage{tasks}
\usepackage{enumerate}
\usepackage{enumitem}
\usepackage{graphicx}

\usepackage[pdftex]{hyperref}
\hypersetup{pdftitle={Física básica},pdfauthor={Luis Carrillo Gutiérrez}}

\fontfamily{cmss}\selectfont
\renewcommand{\familydefault}{\sfdefault}
\pagestyle{empty}

\begin{document}
\paragraph*{Mecánica} -- \\
\begin{itemize}
% 15
\item{Del gráfico, determine el valor de la fuerza entre los bloques. ($g = 10\;m/s^2$)
\begin{tasks}(5)
\task{$\dfrac{2F}{3}$}
\task{$\dfrac{F}{3}$}
\task{$\dfrac{F}{5}$}
\task{$\dfrac{2F}{5}$}
%Ninguna de las anteriores
\end{tasks}
}

% 16
\item{Un bloque de $2\;{Kg}$ se lanza sobre una superficie horizontal rugosa e impacta contra un resorte ($K = 4\;N/cm$). Determine el módulo de su aceleración en el instante en que el resorte está comprimido $3\;cm$ ($g = 10\;m/s^2$)
\begin{tasks}(5)
\task{$4\;m/s^2$}
\task{$6\;m/s^2$}
\task{$8\;m/s^2$}
\task{$10\;m/s^2$}
\task{$12\;m/s^2$}
\end{tasks}
}
\end{itemize}
\end{document}